\subsection{Базовые команды терминала Linux}

\begin{itemize}
	\item \texttt{reset} "--- сброс настроек терминала (например, если сломалась кодировка);
	\item \texttt{clear} или \textit{Ctrl+L} "--- очистка экрана терминала;
	\item \texttt{strings <бинарный файл>} "--- вывод всех текстовых строк, содержащихся в бинарном файле;
	\item \texttt{set} "--- список всех системных переменных;
	\item \texttt{ln myfile myfile.hard} "--- создание жёсткой ссылки (только в пределах файловой системы);
	\item \texttt{ln "="=symbolic ("=s) myfile myfile.sym} "--- создание символической ссылки (может вести на несуществующий файл, аналогично ярлыкам);
	\item \texttt{file <файл>} "--- информация о файле (empty/ASCII string/binary file и т.д.);
	\item \texttt{which <команда>} (например, \texttt{which bash}) "--- информация о расположении бинарного файла команды;
	\item \texttt{mkdir "=p dir1/dir2/dir3} "--- создать иерархическую структуру директорий, все несуществующие директории на пути также будут созданы;
	\item \texttt{rmdir dir1} "--- удалить пустую директорию;
	\item \texttt{rm "=r dir1} "--- удалить директорию вместе со всем её содержимым;
	\item \texttt{ls {>} file.txt} "--- перенаправление вывода в файл (если он уже существовал, то он будет перезаписан);
	\item \texttt{ls {>}{>} file.txt} "--- перенаправление вывода в файл (если он не существовал, то он будет создан, иначе добавление производится в конец существующего файла);
	\item \texttt{history} "--- история команд (можно выполнять навигацию с помощью $\uparrow$ и $\downarrow$);
	\item \texttt{kill "=n <signum> <pid>} "--- послать сигнал с номером \texttt{signum} процессу \texttt{pid};
	\item \texttt{ps "=ef} "--- диспетчер задач.
\end{itemize}

\subsection{Команды для работы с пользователями}

\begin{itemize}
    \item \texttt{w} или \texttt{who} "--- вывод активных пользователей;
    \item \texttt{who am i} или \texttt{whoami} "--- вывод текущего пользователя;
    \item \texttt{id} "--- вывод ID пользователя и его групп;
    \item \texttt{last} "--- вывод даты и времени входов пользователей в систему;
    \item \texttt{passwd <user>} "--- назначение пароля пользователю;
    \item \texttt{su "= <user>} "--- switch user, сменить пользователя;
    \item \texttt{exit} "--- выход из системы;
\end{itemize}

\subsection{Синонимы для директорий}

\begin{itemize}
	\item $\sim$ "--- домашняя директория;
    \item $\sim$ \texttt{<username>} "--- домашняя директория пользователя \texttt{username} (\texttt{/home/username});
	\item . "--- текущая директория;
	\item .. "--- родительская директория.
\end{itemize}

\subsection{Мануалы}

\begin{figure}[H]
	\includegraphics{images/man.jpg}
\end{figure}

Самое важное "--- это ЧИТАТЬ МАНЫ.

\begin{itemize}
	\item \texttt{man <команда>} "--- ман по команде;
	\item \texttt{man "=k <слово>} или \texttt{apropos <слово>} (от appropriate position) "--- все страницы манов, где встретилось слово;
	\item \texttt{<команда> "="=help} или \texttt{<команда> "=h} "--- для простых смертных, краткая справка по команде.
\end{itemize}

\subsection{Работа с правами файлов}

\texttt{ls "=l} выводит подробную информацию обо всех файлах вместе с их флагами.

Пример флага: \texttt{"=rwx"=wx"=w"=}

Подробный разбор флага:
\begin{enumerate}
    \item Первый символ отвечает за то, является ли файл специальным (для директории там будет стоять \texttt{d}, для символической ссылки "--- \texttt{l}, для обычного файла "--- прочерк ("=)).
	\item Следующие три символа (с 2 по 4) отвечают за права владельца.
	\item Следующие три символа (с 5 по 7) отвечают за права группы.
	\item Последние три символа (с 8 по 10) отвечают за права остальных пользователей.
\end{enumerate}

\begin{itemize}
    \item \texttt{r} (Read) "--- право чтения;
    \item \texttt{w} (Write) "--- право записи;
    \item \texttt{x} (eXecute) "--- право исполнения;
	\item "= (прочерк) "--- такого права нет.
\end{itemize}

\texttt{chown <user> <file>} "--- передать права владения файлом пользователю user.

\texttt{chmod u=rwx, g--x, o+rw <file>} "--- дать владельцу все права, у пользователей группы забрать права на исполнение, а остальным пользователям добавить возможность читать файл и писать в него.

Флаг можно закодировать как двоично"=десятичное число. Например, флаг \texttt{rwx"="=x"=wx} будет закодирован как $111_2 001_2 011_2 = 713_{10}$.

\texttt{chmod 777 <file>} "--- дать на файл все права всем пользователям.
